\todo{Lecture 13}

\chapter*{Electromagnetisme}

L'interaction electromagnetique est une des quatres forces fondamentales. Les autres
\begin{enumerate}
\item L'interaction nucleaire forte, portee $\approx 10^{-15}~\mathrm{m}$.
\item L'interaction nucleaire faible, portee $\approx 10^{-17}~\mathrm{m}$.
\item La gravitation, portee illimitee.
\end{enumerate}

L'interaction electromgnetique est:
\begin{itemize}
\item Attractive et repulsive.
\item Responsable de la plupart des phenomenes quotidiens.
\begin{itemize}
\item Electricite.
\item Lumiere.
\item Structure des atoms et des molecules.
\item Gouverne les reactiones chimiques et la biologie.
\item Avec le principe d'exclusion de Pauli (mecanique quantique), la raison pourquoi la terre ne s'ecrasse pas par la gravite, et aussi pourquoi un verre ne tombe pas a travers une table.
\end{itemize}
\item Portee illimitee (mais en pratique les charges positives et negatives tendent a se neutraliser).
\end{itemize}
\chapter{Electrostatique}
\section{Charge electrique}

Il y a deux tpyes de charges appelles positive et negative. Un corps est dit:
\begin{itemize}
\item Neutre, si les charges positives et negatives s'annulent.
\item Chrage, s'il y a un exces d'un type de charge.
\end{itemize}
\subsection{L'unite de la charge}

L'unite de la charge $q$ est le Coulomb $\mathrm{C}=\mathrm{As}$. Dans le Systeme Internationale (SI), le Coulomb est une unite derivee. 

Les septs unites de base du SI sont:
\begin{itemize}
\item $\mathrm{m}$ longueur
\item $\mathrm{kg}$ masse
\item $\mathrm{s}$ temps
\item $\mathrm{A}$ ampere (courant electrique)
\item $\mathrm{K}$ temperature
\item $\mathrm{mol}$ quantite de matiere
\item $\mathrm{cd}$ candela (intensite lumineuse)
\end{itemize}
\begin{example}[Unites derivees]
Le newton est une unite derivee $\mathrm{N}=\mathrm{kgms}^{-2}$.
\end{example}
\begin{remark}
Les definitions des unites de base ont evalue beaucoup historiquement.
\end{remark}

Il existe une charge elementaire $e$:
\begin{itemize}
\item proton: charge $+e$.
\item electron: chagre $-e$.
\end{itemize}
Toutes charges observees sont des multiples de $e$.

\subsection{Conservation de la charge}

La quantite de charge totale est conservee dans toutes les experiences.
\begin{example}[Desintegration de type $\alpha$]
$$^{238}\mathrm{U}\to~^{234}\mathrm{Th}+~^{4}\mathrm{He}$$
Ou $\mathrm{U}$ a 92 protons et 92 electrons, $\mathrm{Th}$ 90 protons et 90 electrons et $\mathrm{He}$ 2 protons et 2 electrons. Seulement la charge totale est conservee, pas la charge positive ou negative.
\end{example}
\begin{example}[Desintegration d'un neutron libre]
$$
n\to p+e^{-}+\bar{\nu}
$$
ou $\bar{\nu}$ est neutre.
\end{example}
\begin{example}[Creation des paires d'electron et positron par un photon $\gamma$]
$$
\gamma \to e^{-}+e^{ + }
$$
\end{example}

\section{Loi de Coulomb}

Soient deux charges ponctuelles $q_{1}$ et $q_{2}$ aux positions $\boldsymbol{r}_{1}$ et $\boldsymbol{r}_{2}$ respectivement et $\boldsymbol{r}_{12}$ un vecteur entre ces deux charges. La force exercee par $q_{1}$ sur $q_{2}$ est 
$$
\boldsymbol{F}_{1\to 2}=\frac{1}{4\pi\varepsilon _{0}} \frac{q_{1}q_{2}}{\lVert \boldsymbol{r}_{12} \rVert ^{2}} \frac{\boldsymbol{r}_{12}}{\lVert \boldsymbol{r}_{12} \rVert }
$$
ou de maniere equivalente
$$
\boldsymbol{F}_{1\to 2}=\frac{1}{4\pi\varepsilon_{0}} \frac{q_{1}q_{2}}{\lVert \boldsymbol{r}_{2}-\boldsymbol{r}_{1} \rVert ^{2}} \frac{\boldsymbol{r}_{2}-\boldsymbol{r}_{1}}{\lVert \boldsymbol{r}_{2}-\boldsymbol{r}_{1} \rVert }
$$
ou $\varepsilon _{0}=8.85\cdot 10^{-12}~\mathrm{Fm}^{-1}$ ou $F=\mathrm{A}^{2}\mathrm{s}^{4}\mathrm{kg}^{-1}\mathrm{m}^{-2}$ est un Farad.

Par la troisieme loi de Newton on a
$$
\boldsymbol{F}_{1\to 2} =-\boldsymbol{F}_{2 \to 1}
$$
\subsection{Principe de superposition}

Soient trois charges $q_{1}$, $q_{2}$ et $q_{3}$ ponctuelles. La force $\boldsymbol{F}_{1 \to 3}$ est la force de $q_{1}$ sur $q_{3}$ en abscense de $q_{2}$ et $\boldsymbol{F}_{2\to 3}$ est la force de $q_{2}$ sur $q_{3}$ en abscence de $q_{1}$. La force en presence de $q_{1}$ et $q_{2}$ est donc
$$
\boldsymbol{F}_{\mathrm{res}}=\boldsymbol{F}_{1\to 3} +\boldsymbol{F}_{2\to 3}
$$

Force de $n$ charges ponctuelles $q_{i}$ a $\boldsymbol{r}_{i}$ sur une charge $q$ a $\boldsymbol{r}$:
$$
\boldsymbol{F}_{\mathrm{res}}=\sum_{i=1}^{n} \frac{1}{4\pi\varepsilon_{0}} \frac{qq_{i}}{\lVert \boldsymbol{r} -\boldsymbol{r}_{i}\rVert^{2}} \frac{\boldsymbol{r}-\boldsymbol{r}_{i}}{\lVert \boldsymbol{r}-\boldsymbol{r}_{i} \rVert }
$$
\section{Champ electrique}

Soit une charge $q_{1}$ en $\boldsymbol{r}_{1}$. A une position $\boldsymbol{r}$ de l'espace, une charge $q$ va subir la force
$$
\boldsymbol{F}_{1\to q} =\underbrace{ \frac{1}{4\pi\varepsilon_{0}} \frac{q_{1}}{\lVert \boldsymbol{r}-\boldsymbol{r}_{1} \rVert ^{2}} \frac{\boldsymbol{r}-\boldsymbol{r}_{1}}{\lVert \boldsymbol{r}-\boldsymbol{r}_{1} \rVert } }_{ \text{Champ vectoriel independent de }q } \cdot q
$$
On appelle ce champ le champ electrique $\boldsymbol{E}$ genere par $q_{1}$.

Le principe de superposition implique que le champ electrique genere par $n$ charges $q_{i}$ a poistion $\boldsymbol{r}_{i}$ :
$$
\boldsymbol{E}(\boldsymbol{r})=\sum_{i=1}^{n} \frac{1}{4\pi\varepsilon_{0}} \frac{q_{i}}{\lVert \boldsymbol{r}-\boldsymbol{r}_{i} \rVert ^{2}} \frac{\boldsymbol{r}-\boldsymbol{r}_{i}}{\lVert \boldsymbol{r}-\boldsymbol{r}_{i} \rVert }
$$
donc 
$$
\boldsymbol{E}=\frac{\boldsymbol{F}}{q}
$$
le champ electrique est dans la direction de la force ressentie par une charge positive.

Le champ electrique est mesure en volt-metre
$$
[\boldsymbol{E}]=\frac{\mathrm{N}}{\mathrm{C}}=\mathrm{kgms}^{-2}\mathrm{C}=\mathrm{kgms}^{-2}\mathrm{A}^{-1}\mathrm{s}^{-1}=\frac{\mathrm{V}}{\mathrm{m}}
$$